%!LW recipe=latexmk (xelatex)
\documentclass[12pt, a4paper, oneside]{ctexbook}
\usepackage{amsmath, amsthm, amssymb, bm, graphicx, hyperref, mathrsfs, epigraph, footnote, cleveref, esint, wrapfig, caption, appendix, microtype}
\usepackage[backend=biber, style=gb7714-2015]{biblatex}

\usepackage{import}
\usepackage{xifthen}
\usepackage{pdfpages}
\usepackage{transparent}

\newcommand{\incfig}[1]{%
    \def\svgwidth{0.5\columnwidth}
    \import{./figures/}{#1.pdf_tex}
}

% \addbibresource{myrefs.bib}
\addbibresource{recommend.bib}

\title{{\Huge{\textbf{电动力学笔记}}}}
\author{ZexHexoft}
\date{\today}
\linespread{1.5}
\newtheorem{theorem}{定理}[section]
\newtheorem{definition}[theorem]{定义}
\newtheorem{lemma}[theorem]{引理}
\newtheorem{corollary}[theorem]{推论}
\newtheorem{example}[theorem]{例}
\newtheorem{proposition}[theorem]{命题}
\newenvironment{problem}{\begin{shaded}\stepcounter{problemname}\par\noindent\textbf{题目\arabic{problemname}. }}{\end{shaded}\par}
\newenvironment{solution}{\par\noindent\textbf{解答. }}{\par}
\newenvironment{note}{\par\noindent\textbf{题目\arabic{problemname}的注记. }}{\par}

\setmainfont{Times New Roman}

% \def\rcurs{{\mbox{$\resizebox{.08in}{.08in}{\includegraphics{./script_r/ScriptR}}$}}}
% \def\brcurs{{\mbox{$\resizebox{.08in}{.08in}{\includegraphics{./script_r/BoldR}}$}}}
% \def\hrcurs{{\mbox{$\hat \brcurs$}}}

\newcommand{\rcurs}{{\mbox{$\resizebox{.08in}{.08in}{\includegraphics{./script_r/ScriptR}}$}}}
\newcommand{\brcurs}{{\mbox{$\resizebox{.08in}{.08in}{\includegraphics{./script_r/BoldR}}$}}}
\newcommand{\hrcurs}{{\mbox{$\hat \brcurs$}}}

\renewcommand{\vec}[1]{\mathbf{#1}}
\newcommand{\uvec}[1]{\hat{\mathbf{#1}}}
\newcommand{\rmd}{\mathrm{d}}

\crefname{figure}{图}{图}
\Crefname{figure}{图}{图}
\crefname{table}{表}{表}
\Crefname{table}{表}{表}
\crefname{equation}{式}{式}
\Crefname{equation}{式}{式}
\crefname{theorem}{定理}{定理}
\Crefname{theorem}{定理}{定理}
\crefname{section}{}{}
\Crefname{section}{}{}
\crefname{chapter}{}{}
\Crefname{chapter}{}{}

\crefformat{section}{\textbf{#2#1#3}部分}
\crefformat{chapter}{#2第#1章#3}
\crefformat{equation}{式#2(#1)#3} % 这里的几个格式非常有趣
% \crefformat{equation}{#2(#1)#3}
% \crefformat{equation}{#2式(#1)#3}

\crefformat{figure}{#2图#1#3}

\newcommand{\crefpairconjunction}{ 和 }
\newcommand{\creflastconjunction}{ 和 }
\newcommand{\crefmiddleconjunction}{、}
\newcommand{\crefrangeconjunction}{ 至 }

\renewcommand{\appendixname}{附录}
\renewcommand{\appendixpagename}{附录}
\renewcommand{\appendixtocname}{附录}

\hypersetup{
    colorlinks=true,
    linkcolor=blue,
    urlcolor=blue,
    citecolor=red
}

\begin{document}
\pagenumbering{Alph}
\maketitle

\pagenumbering{roman}
\setcounter{page}{1}

\begin{center}
    \Huge\textbf{前言}
\end{center}~\

这是我本人在学习电动力学时的一些笔记和想法.

希望各位能从中汲取一些有用的想法.

末尾有一份我本人的推荐书单，有兴趣可以去阅读(书单顺序不代表阅读顺序).

本文目前在\href{https://github.com/ZexHexoft/ElectroDynamicsNotes}{Github}上发布.

注: 目前这篇文档尚未完成，且限于作者水平，难免存在纰漏. 如有问题或建议，烦请各位在 \href{https://github.com/ZexHexoft/ElectroDynamicsNotes/issues}{issue} 中指正.

~\\
\begin{flushright}
    \begin{tabular}{c}
        ZexHexoft\\
        2026级复旦大学物理学系本科生\\
        \today
    \end{tabular}
\end{flushright}


\newpage
\pagenumbering{Roman}
\setcounter{page}{1}
\tableofcontents
\newpage

\setcounter{chapter}{-1}
\setcounter{page}{1}
\pagenumbering{arabic}



% !TEX root = ../notes.tex

\chapter{符号定义}

% \epigraph{数学是描述世界的语言，但就像世界上大多数语言一样，总有人能用一些合理的符号组成令人迷惑的句子.}{——ZexHexoft\\于2026.9.8晚上23:22}

书中主要采用与David J.Griffiths所著的《电动力学导论(Introduction To Electrodynamics)》类似的记号(但有细微的不同).

具体而言：
\begin{itemize}
    \item 在这里，我们使用诸如$\mathbf{r}$,$\mathbf{v}$,$\mathbf{a}$的符号表示矢量
    \item 而$r,v,a$则对应他们的"大小"
    \item $\hat{\mathbf{r}}$,  $\hat{\mathbf{v}}$,  $\hat{\mathbf{a}}$表示同方向的单位向量
    \item 特殊的，我们一般用$\mathbf{r}$表示位置矢量(由原点指向该位置), 用$\mathbf{v}$表示速度
    \item 更特殊的，我们使用$\brcurs$ 表示"间隔矢量"(即$\brcurs=\mathbf{r}-\mathbf{r}'$, 从源点指向场点), 用$\rcurs$表示其大小
    \item 对于某个量(例如$r,\mathbf{v}$)对时间的一阶导数，我们一般写成$\dot{r},\dot{\mathbf{v}}$的形式, 求几阶导就加几个点, 甚至可以$\overset{........}{y}$ 
    \item 在积分运算中，我们会使用,$\mathrm{d}\mathbf{a}$或者$\mathrm{d}\mathbf{S}$表示面积微元, 用$\mathrm{d}\tau$或者$\mathrm{d}^3\vec x$表示体积微元.
    \item 另外，我们有时会使用爱因斯坦求和规则，阅读时请注意.
\end{itemize}

有人可能会觉得这里的符号有点奇怪(也许确实是这样), 但谁让我读的第一本电动力学教材是Griffiths的, 而你又选择了读这本笔记呢?
\footnote{当然，这里的所有记号规范可能只是我的一厢情愿，实际请以文意为准.}



\chapter{电动力学预备班}
\epigraph{你应该已经学过微积分和矢量分析了, 对吧?}{——ZexHexoft}
这部分主讲常用矢量分析和一些相关的定理及公式.\footnote{事实上我这里非常推荐阅读《电磁学(拓展篇)》中的专题15, 接下来也会引用部分内容}

\section{常用的矢量恒等式}

\begin{theorem}[标量三重积]
    \begin{equation}
        \vec A\cdot(\vec B\times\vec C)=\vec B\cdot(\vec C\times\vec A)=\vec C\cdot(\vec A\times\vec B)
    \end{equation}
\end{theorem}

\begin{theorem}[矢量三重积(BAC-CAB规则\footnote{英文中与BACK-CAB同音, 即后面的出租车})] \label{BAC-CAB}
        \begin{equation}
            \vec A\times(\vec B\times\vec C)=\vec B(\vec A\cdot\vec C)-\vec C(\vec A\cdot\vec B)
        \end{equation}
\end{theorem}

\section{矢量的微分}

\subsection{一阶微分}

\begin{definition}[$\nabla$算子]
    \begin{equation}
        \nabla := \hat{\mathbf{x}}\frac{\partial}{\partial x}+\hat{\mathbf{y}}\frac{\partial}{\partial y}+\hat{\mathbf{z}}\frac{\partial}{\partial z}
    \end{equation}
\end{definition}

\begin{theorem}[梯度]
    \begin{equation}
        \nabla T := \frac{\partial T}{\partial x}\hat{\mathbf{x}}+\frac{\partial T}{\partial y}\hat{\mathbf{y}}+\frac{\partial T}{\partial z}\hat{\mathbf{z}}
    \end{equation}
\end{theorem}

\begin{theorem}[散度]
    \begin{equation}
        \nabla \cdot \mathbf{v}=\frac{\partial v_x}{\partial x}+\frac{\partial v_y}{\partial y}+\frac{\partial v_z}{\partial z}
    \end{equation}
\end{theorem}

在\cref{sec:integrationOfVector}中，我们会看到他们名字的意义.

\begin{theorem}[旋度]\footnote{都学到这里了，你会计算这个的，对吧? 那我就不展开了.}
    \begin{equation}
        \nabla \times \mathbf{v}=\left| \begin{matrix}
            \hat{\mathbf{x}}&		\hat{\mathbf{y}}&		\hat{\mathbf{z}}\\
            \frac{\partial}{\partial x}&		\frac{\partial}{\partial y}&		\frac{\partial}{\partial z}\\
            v_x&		v_y&		v_z\\
        \end{matrix} \right|
    \end{equation}
\end{theorem}

\begin{theorem}[微分的积规则]
    \begin{equation}
        \nabla(fg) = f\nabla g+g\nabla f
    \end{equation}
    \begin{equation}
        \nabla(\vec A \cdot \vec B) = \vec A \times (\nabla \times \vec B) + \vec B \times (\nabla \times \vec A) + (\vec A \cdot \nabla)\vec B+ (\vec B \cdot \nabla)\vec A \label{eq:example2}
    \end{equation}
    \begin{equation}
        \nabla \cdot \left( f\mathbf{A} \right) =f\left( \nabla \cdot \mathbf{A} \right) +\mathbf{A}\cdot \nabla f
    \end{equation}
    \begin{equation}
        \nabla \cdot \left( \mathbf{A}\times \mathbf{B} \right) =\mathbf{B}\cdot \left( \nabla \times \mathbf{A} \right) -\mathbf{A}\cdot \left( \nabla \times \mathbf{B} \right) 
    \end{equation}
    \begin{equation}
        \nabla \times \left( f\mathbf{A} \right) =f\left( \nabla \times \mathbf{A} \right) -\mathbf{A}\times \left( \nabla f \right) 
    \end{equation}
    \begin{equation}
        \nabla \times \left( \mathbf{A}\times \mathbf{B} \right) =\left( \mathbf{B}\cdot \nabla \right) \mathbf{A}-\left( \mathbf{A}\cdot \nabla \right) \mathbf{B}+\mathbf{A}\left( \nabla \cdot \mathbf{B} \right) -\mathbf{B}\left( \nabla \cdot \mathbf{A} \right) \label{eq:example1}
    \end{equation}
\end{theorem}

说实话，这些等式又臭又长，很多时候难以记住.

我在这里提供一种摘自《电磁学(拓展篇)》中的专题15中选读15-10的一种记忆方法供大家参考.

首先简记一种符号$\nabla_\vec{A}$(可以与一元微积分中的$\frac{\mathrm{d}f}{\mathrm{d}x}\frac{\partial}{\partial f},\frac{\mathrm{d}g}{\mathrm{d}x}\frac{\partial}{\partial g}$等对应)

\begin{definition}
    $\nabla_\vec{A}$表示只作用于向量$\vec{A}$的算符.
\end{definition}

这里以\cref{eq:example1}为例子

\begin{example}
    \begin{align*}
        \nabla \times \left( \mathbf{A}\times \mathbf{B} \right) &=\nabla _{\mathbf{A}}\times \left( \mathbf{A}\times \mathbf{B} \right) +\nabla _{\mathbf{B}}\times \left( \mathbf{A}\times \mathbf{B} \right) 
        \\
        &=\mathbf{A}\left( \nabla _{\mathbf{A}}\cdot \mathbf{B} \right) -\mathbf{B}\left( \nabla _{\mathbf{A}}\cdot \mathbf{A} \right) +\mathbf{A}\left( \nabla _{\mathbf{B}}\cdot \mathbf{B} \right) -\mathbf{B}\left( \nabla _{\mathbf{B}}\cdot \mathbf{A} \right)
        \\
        &=\left( \mathbf{B}\cdot \nabla _{\mathbf{A}} \right) \mathbf{A}-\mathbf{B}\left( \nabla _{\mathbf{A}}\cdot \mathbf{A} \right) +\mathbf{A}\left( \nabla _{\mathbf{B}}\cdot \mathbf{B} \right) -\left( \mathbf{A}\cdot \nabla _{\mathbf{B}} \right) \mathbf{B}
        \\
        &=\left( \mathbf{B}\cdot \nabla \right) \mathbf{A}-\mathbf{B}\left( \nabla \cdot \mathbf{A} \right) +\mathbf{A}\left( \nabla \cdot \mathbf{B} \right) -\left( \mathbf{A}\cdot \nabla \right) \mathbf{B}
    \end{align*}
\end{example}

再举一个\cref{eq:example2}的例子(注意对\cref{BAC-CAB}的逆运用)

\begin{example}
    \begin{align*}
        \nabla \left( \mathbf{A}\cdot \mathbf{B} \right) &=\nabla _{\mathbf{A}}\left( \mathbf{A}\cdot \mathbf{B} \right) +\nabla _{\mathbf{B}}\left( \mathbf{A}\cdot \mathbf{B} \right)
        \\
        &=\mathbf{B}\times \left( \nabla _{\mathbf{A}}\times \mathbf{A} \right) +\mathbf{A}\left( \mathbf{B}\cdot \nabla _{\mathbf{A}} \right) +\mathbf{A}\times \left( \nabla _{\mathbf{B}}\times \mathbf{B} \right) +\mathbf{B}\left( \mathbf{A}\cdot \nabla _{\mathbf{B}} \right) 
        \\
        &=\mathbf{B}\times \left( \nabla \times \mathbf{A} \right) +\left( \mathbf{B}\cdot \nabla \right) \mathbf{A}+\mathbf{A}\times \left( \nabla \times \mathbf{B} \right) +\left( \mathbf{A}\cdot \nabla \right) \mathbf{B}
    \end{align*}
\end{example}

\subsection{常用的二阶微分}

\begin{definition}[拉普拉斯算符]\footnote{似乎有点反直觉的是，拉普拉斯算符既可以作用在标量上，也可以作用在矢量上.}
    \begin{equation}
        \nabla^2 := \nabla \cdot \nabla = \frac{\partial ^2}{\partial x^2}+\frac{\partial ^2}{\partial y^2}+\frac{\partial ^2}{\partial z^2}
    \end{equation}
\end{definition}

\begin{theorem}
    \begin{gather}
        \nabla \times \left( \nabla T \right) =0\\
        \nabla \cdot \left( \nabla \times \mathbf{v} \right) =0\\
        \nabla \times \left( \nabla \times \mathbf{v} \right) =\nabla \left( \nabla \cdot \mathbf{v} \right) -\nabla ^2\mathbf{v}
    \end{gather}
\end{theorem}


\section{矢量的积分}\label{sec:integrationOfVector}

\begin{theorem}[梯度定理]
    \begin{equation}
        \int_a^b{\left( \nabla T \right) \cdot \mathrm{d}\mathbf{l}}=T\left( b \right) -T\left( a \right) 
    \end{equation}
\end{theorem}

\begin{theorem}[散度定理]
    \begin{equation}
        \int_V{\left( \nabla \cdot \mathbf{v} \right) \mathrm{d}\tau}=\oint_S{\mathbf{v}\cdot \mathrm{d}\mathbf{a}}
    \end{equation}
\end{theorem}

\begin{theorem}[旋度定理/斯托克斯定理]\footnote{这里的$\partial S$表示$S$的边界.}\footnote{之所以在积分号上用逆时针的圆圈，是为了强调右手规则.}
    \begin{equation}
        \int_S{\left( \nabla \times \mathbf{v} \right) \cdot \mathrm{d}\mathbf{a}}=\ointctrclockwise_{\partial S}{\mathbf{v}\cdot \mathrm{d}\mathbf{l}}
    \end{equation}
\end{theorem}
\chapter{静电学}

% \epigraph{我知道这是静电学的内容，但我就是要放在这里，你气不气}{——ZexHexoft\\于码字码疯时的话语}

\section{静电场}

\subsection{真空中的静电场}

\begin{definition}[静电场]
    \begin{equation}
        \mathbf{E}=\frac{1}{4\pi \varepsilon _0}\int_V{\frac{\rho \left( \mathbf{r}' \right)}{\rcurs^2}\hrcurs\mathrm{d}\tau '}
    \end{equation}
\end{definition}

\begin{theorem}[高斯定理]
    \begin{gather}
        \oint_S{\mathbf{E}\cdot \mathrm{d}\mathbf{a}}=\frac{Q}{\varepsilon _0}\\
        \nabla \cdot \mathbf{E}=\frac{\rho}{\varepsilon _0}
    \end{gather}
\end{theorem}

\begin{definition}[静电势]
    \begin{equation}
        V(\vec r):=-\int_O^r{\mathbf{E}\cdot \mathrm{d}\mathbf{l}}
    \end{equation}
\end{definition}

\begin{corollary}
    \begin{equation}
        \vec E = - \nabla V
    \end{equation}
\end{corollary}

\begin{definition}[泊松方程]
    \begin{equation}
        \nabla^2 V = - \frac{\rho}{\varepsilon_0} 
    \end{equation}
\end{definition}

对无电荷区域，退化为

\begin{definition}[拉普拉斯方程]
    \begin{equation}
        \nabla^2 V = 0
    \end{equation}
\end{definition}

\begin{theorem}[边界条件]\label{boundary_statement}
    \begin{gather}
        \mathbf{E}_{\text{上}}-\mathbf{E}_{\text{下}}=\frac{\sigma}{\varepsilon _0}\hat{\mathbf{n}}\\
        \frac{\partial V_{\text{上}}}{\partial n}-\frac{\partial V_{\text{下}}}{\partial n}=-\frac{\sigma}{\varepsilon _0}
    \end{gather}
\end{theorem}

我们熟知分立电荷的能量

\begin{theorem}
    \begin{equation}
        W=\frac{1}{2}\sum_{i=1}^n{q_iV\left( \mathbf{r}_i \right)}
    \end{equation}
\end{theorem}

接下来我们计算连续电荷分布的能量\footnote{注意\cref{eq:diff_rule3}的逆运用}

\begin{theorem}
    \begin{align}
        W&=\frac{1}{2}\int{\rho V\mathrm{d}\tau}
        \\
        &=\frac{\varepsilon _0}{2}\int{\left( \nabla \cdot \mathbf{E} \right) V\mathrm{d}\tau}\nonumber
        \\
        &=\frac{\varepsilon _0}{2}\int{\left[ \nabla \cdot \left( V\mathbf{E} \right) -\mathbf{E}\cdot \nabla V \right] \mathrm{d}\tau}\nonumber
        \\
        &=\frac{\varepsilon _0}{2}\left( \int_V{E^2\mathrm{d}\tau}+\oint_S{V\mathbf{E}\cdot \mathrm{d}\mathbf{a}} \right) \nonumber
        \\
        &=\frac{\varepsilon _0}{2}\int_V{E^2\mathrm{d}\tau}
    \end{align}
\end{theorem}

\subsection{介质中的静电场}

考虑线性电介质:

\begin{theorem}[束缚电荷]
    \begin{gather}
        Q_b:=\int{\mathrm{d}^3\mathbf{x}\rho _b}=-\oint{\mathbf{P}\cdot \mathrm{d}\mathbf{a}}\\
        \rho_b = -\nabla \cdot \vec P\\
        \sigma_b = \hat{\vec n}\cdot \vec P
    \end{gather}
\end{theorem}

由连续性方程

$$
\frac{\partial \rho_b}{\partial t}+\nabla \cdot \vec J_b = 0
$$

我们有如下推论: 

\begin{corollary}[束缚电流密度]
    \begin{equation}
        \vec J_b = \frac{\partial \vec P}{\partial t}
    \end{equation}
\end{corollary}

\section{导体}

我们熟知导体有以下性质：

\begin{itemize}
    \item 内部$\vec E = 0$
    \item 内部$\rho = 0$
    \item 净电荷只分布在表面
    \item 是一个等势体
    \item 电场垂直于导体表面
\end{itemize}





\begin{refsection}[recommend.bib]
  \nocite{*}
  \printbibliography[title=推荐书单]
\end{refsection}

\end{document}